\documentclass[11pt]{article}
\usepackage[margin=1in]{geometry}
\usepackage[T1]{fontenc}
\usepackage[utf8]{inputenc}
\usepackage{lmodern,microtype,amsmath,amssymb,amsthm,mathtools,booktabs}
\usepackage[colorlinks=true,linkcolor=blue,citecolor=blue,urlcolor=blue]{hyperref}
\newtheorem{theorem}{Theorem}
\newtheorem{proposition}[theorem]{Proposition}
\newtheorem{lemma}[theorem]{Lemma}
\newtheorem{corollary}[theorem]{Corollary}
\newtheorem{openproblem}[theorem]{Open atomic proposition}
\theoremstyle{remark}\newtheorem{remark}[theorem]{Remark}
\newcommand{\E}{\mathbb E}\newcommand{\R}{\mathbb R}
\newcommand{\F}{\mathbb F}\newcommand{\one}{\mathbf 1}
\newcommand{\sech}{\operatorname{sech}}
\newcommand{\Var}{\operatorname{Var}}
\title{Superpolynomial online mistakes for isotropic Bayesian tanh networks}
\author{ChatGPT-6 Astra \and Kaarel H\"anni \and Dmitry Vaintrob\thanks{Authors are listed in alphabetical order.}}
\date{November 2026}
\hypersetup{pdftitle={Superpolynomial online mistakes for isotropic Bayesian tanh networks},pdfauthor={ChatGPT-6 Astra; Kaarel Hanni; Dmitry Vaintrob}}
\begin{document}\maketitle
\begin{abstract}
We construct small Boolean targets and fixed nonrepeating input
sequences on which exact Bayesian learning with an ordinary tanh
network makes superpolynomially many online mistakes. The network has
one hidden layer, a linear readout, and independent Gaussian weights
whose covariance is isotropic within each layer. The prior is
conditioned from the start on an absolute output margin at every
Boolean input, and training requires exact fit of all revealed labels.
At every sufficiently large fixed learner width, the posterior
probability of the correct label tends to zero on each designated
round as the Gaussian scale tends to zero.

The targets have one-hidden-layer tanh teachers of width $C$, margin
one, and coefficients bounded by $C$. The construction forces
$\exp(\Omega(C^{1/8}))$ designated mistakes. Its calibration prefix
and sufficient learner width can be large. This is an adversarial-order
result, rather than an iid sample-complexity lower bound.

The proof first determines the exact optimal affine neuron for a
single sufficiently high odd parity. A quantitative support estimate
and Fourier localization extend the sharp norm to a transversal code.
The optimized scalar costs have positive inverse-degree curvature;
their cheapest feasible scores therefore include a fixed negative
overshoot. An explicit dual classifies every minimum, a cycle code
orders superpolynomially many competing costs, and integer neuron
replication transfers their gaps to the full Gaussian posterior.
\end{abstract}

\section{Model and result}

A small circuit may be easy to learn from independent samples while
remaining difficult along an adversarial order of inputs. We study
that distinction for exact Bayesian prediction with the standard
activation $\tanh$. The theorem concerns the posterior distribution,
including all tied or nearly tied parameter realizations.

On $X=\{-1,1\}^d$, append the constant input and write $\bar x=(1,x)$.
Consider a one-hidden-layer network
\[
 F(x)=\sum_{j=1}^N v_j\tanh(u_j\cdot\bar x),\qquad
 (u_j)_i\sim\mathcal N(0,N\sigma^2/d),\qquad
 v_j\sim\mathcal N(0,\sigma^2/N),
\]
with independent coordinates. The prior is first conditioned on
$|F(x)|>1/2$ for every input; training then conditions on the exact
observed labels. The strong-prior limit is $\sigma\downarrow0$ at each
fixed width $N$. Equivalently, at that fixed width set
$t=\sqrt N\,\sigma$: the variances become $t^2/d$ in the incoming
layer and $t^2/N^2$ in the readout. Thus this is the usual strong-prior
limit, with isotropy within each layer preserved. Its quadratic energy is
\[
 K_N=\frac dN\sum_j\|u_j\|_2^2+N\sum_jv_j^2.
\]
For a function on the finite cube, define the tanh variation cost
\begin{equation}\label{eq:variation}
 V(F)=\inf\left\{\sum_j|c_j|\|a_j\|_2:
 F(x)=\sum_jc_j\tanh(a_j\cdot\bar x)\text{ for all }x\right\}.
\end{equation}
The infimum is over finite representations and includes the incoming
bias coordinate in the norm. The arithmetic--geometric mean inequality
always gives
\begin{equation}\label{eq:energy}
 K_N\ge2\sqrt d\sum_j|v_j|\|u_j\|_2\ge2\sqrt d\,V(F).
\end{equation}
Unlike a homogeneous activation, tanh does not permit arbitrary
rescaling of an incoming vector while preserving the neuron output up
to a scalar. Equality in \eqref{eq:energy} at a specified finite width
therefore requires a separate argument.

\begin{theorem}[Adversarial online mistakes]\label{thm:main-online}
For every sufficiently large integer $k$, put
$m=3k(k+5)/2+3$. There are a dimension $d_m=\Theta(m^3)$,
a Boolean target $f_m$ on the full cube, and a fixed nonrepeating
input sequence with $T_m=\exp(\Omega(\sqrt m))$ designated rounds
such that the following hold.
The target has a one-hidden-layer tanh score of margin one and width
$C_m=O(m^4)$, with every coefficient bounded by $C_m$. It also
has a bounded-fan-in Boolean circuit with $O(d_m)=O(m^3)$ gates.
There is a finite $N_0(m)$ such that, at every fixed $N\ge N_0(m)$,
the posterior probability of the correct label tends to zero at each
designated round as $\sigma\downarrow0$.
Consequently the sequence forces $\exp(\Omega(C_m^{1/8}))$
posterior-majority mistakes in the strong-prior limit, and the same
number of expected mistakes under posterior sampling.
\end{theorem}

\begin{corollary}[Finite Gaussian scales]\label{cor:finite-scale}
Fix a construction, any $N\ge N_0(m)$, and $\eta>0$. For every
sufficiently small positive $\sigma$, posterior-majority prediction
is wrong at all $T_m$ designated rounds, and the correct-label
posterior probability is below $1/3$ at every such round.
Under fresh posterior sampling, all $T_m$ designated predictions are
wrong simultaneously with probability at least $1-\eta$.
The input sequence can be extended to a permutation of the full cube.
\end{corollary}
\begin{proof}
There are finitely many designated rounds. Choose $\sigma$ so that
the correct-label probability at each is below
$\min(1/3,\eta/T_m)$. Majority and the confidence-threshold conclusion
follow, while a union bound proves the sampling statement. Append
all inputs not yet revealed to extend the sequence without repetition.
\end{proof}

The order of limits is part of the statement: first choose any
sufficiently large finite learner width, then send the Gaussian scale
to zero. No distribution-shift or typical-order claim is assumed.
The input sequence has a potentially large calibration prefix and
is independent of the learner width and Gaussian scale. The sufficiently
large parity-degree threshold and the width threshold are existential;
this proof does not give practical numerical values for them. The
arguments in this manuscript have not been formalized in Lean.

All Fourier coefficients use the uniform input law. For a set $A$,
write $\chi_A(x)=\prod_{i\in A}x_i$ and
$\widehat F(A)=\E[\chi_A F]$. A sharp inequality of the form
\[
 |\widehat{\tanh(u\cdot\bar x)}(A)|\le L_A\|u\|_2
\]
gives the lower bound $V(\chi_A)\ge L_A^{-1}$. Signed averaging over
input flips gives a matching parity representation whenever a neuron
attains the sharp correlation ratio. Controlling a sum of several
parity coefficients is stronger and is needed for the online construction.

\section{Balanced odd-parity atoms}

Let $n=2r+1$ be odd and let $X_1,\ldots,X_n$ be independent uniform signs.
Put $S=\sum_iX_i$, $\chi=\prod_iX_i$, and
\[
 q_n(a)=(-1)^r\E[\chi\tanh(aS)],\qquad
 \kappa_n=2^{1-n}\binom{n-1}{(n-1)/2}.
\]
The constant $\kappa_n$ is the absolute top Fourier coefficient of
majority on $n$ bits. The balanced representation cost is
\begin{equation}\label{eq:balanced-cost}
 w_n=\inf_{a>0}\frac{a\sqrt n}{q_n(a)}.
\end{equation}
The positivity needed in this definition follows from the next formula.

\begin{lemma}[A positive diagonal series]\label{lem:positive-series}
For every odd $n\ge1$ and $a>0$, with $c_\ell=(\ell+1/2)\pi$,
\begin{equation}\label{eq:positive-series}
 q_n(a)=2n!a^n\sum_{\ell\ge0}
 \prod_{j=0}^{(n-1)/2}\bigl(c_\ell^2+(2j+1)^2a^2\bigr)^{-1}>0.
\end{equation}
\end{lemma}
\begin{proof}
The sine transform of tanh is
\[
 \tanh t=\int_0^\infty\frac{\sin(t\omega)}{\sinh(\pi\omega/2)}\,d\omega.
\]
Averaging the parity character gives
\[
 q_n(a)=\int_0^\infty
 \frac{\sin^n(a\omega)}{\sinh(\pi\omega/2)}\,d\omega.
\]
Expand $1/\sinh(\pi\omega/2)=2\sum_{\ell\ge0}e^{-c_\ell\omega}$.
Absolute integrability of $|\sin(a\omega)|^n/\sinh(\pi\omega/2)$
justifies the exchange. Repeated integration by parts gives
\[
 \int_0^\infty e^{-c\omega}\sin^n(a\omega)\,d\omega
 =\frac{n!a^n}{\prod_{j=0}^{(n-1)/2}(c^2+(2j+1)^2a^2)}.
\]
This proves the formula and positivity.
\end{proof}

\begin{theorem}[Unique balanced gain and full local optimality]
\label{thm:local}
For every odd $n\ge5$, the ratio $q_n(a)/a$ has a unique global maximum
at a gain $a_n\in(1,2)$. Moreover
\[
 J_n(b,u)=\frac{(-1)^{(n-1)/2}\E[\chi\tanh(b+u\cdot X)]}
                   {\sqrt{b^2+\|u\|_2^2}}
\]
has a strict local maximum at $(b,u)=(0,a_n\one)$. Its Hessian there is
negative definite. The absolute Fourier-correlation ratio equals $J_n$
in a neighborhood of that point.
\end{theorem}

\begin{proof}
For any odd function $h$, pairing opposite Hamming levels yields
\begin{equation}\label{eq:central-sum}
 (-1)^r\E[\chi h(aS)]
 =2^{1-n}\sum_{j=0}^r(-1)^j\binom n{r-j}h(a(2j+1)).
\end{equation}
Let $p(t)=-\tanh''t=2\tanh t\,\sech^2t$. It is positive and strictly
decreasing for $t\ge1$. Consequently
\begin{equation}\label{eq:Dnegative}
 D_n(a):=(-1)^r\E[\chi\tanh''(aS)]<0
 \qquad(a\ge1)
\end{equation}
for every odd $n\ge1$: the positive terms in the corresponding
alternating sum strictly decrease.

The same argument proves $q_n''(a)<0$ on $a\ge1$. The extra factor in
its terms is $(2j+1)^2$. For $t'=t+2a$, $t\ge a\ge1$,
\[
 \frac{p(t')}{p(t)}
 \le e^{-4a}\frac{(1+e^{-2a})^3}{1-e^{-2a}}.
\]
The ratio of the squared integer factors is at most nine. Since $e^2>7$,
\[
 9e^{-4a}\frac{(1+e^{-2a})^3}{1-e^{-2a}}
 <\frac9{49}\frac{(8/7)^3}{6/7}<1.
\]
The binomial factors decrease as well, giving the claimed strict sign.

For $0<a\le1$, every summand of $q_n(a)/a$ in
\eqref{eq:positive-series} has positive derivative. Its logarithmic
derivative multiplied by $a$ is
\[
 -2+2\sum_{j=0}^r
       \frac{c_\ell^2}{c_\ell^2+(2j+1)^2a^2}.
\]
Using $n\ge5$ and $c_\ell^2\ge\pi^2/4>12/5$, the sum is greater than
\[
 \frac{12}{17}+\frac{12}{57}+\frac{12}{137}
 =\frac{44428}{44251}>1.
\]
The differentiated series is locally uniformly convergent on $a>0$,
so $(q_n(a)/a)'>0$ throughout $(0,1]$.

Put $H_n(a)=a q_n'(a)-q_n(a)$. Its derivative is $a q_n''(a)<0$
for $a\ge1$. Define
\[
 g(t)=t\sech^2t+1-\tanh t,
 \qquad g'(t)=-2t\sech^2t\tanh t<0.
\]
Subtracting the majority coefficient in \eqref{eq:central-sum}, and
using alternating-sum bounds, gives
\[
 H_n(1)\ge-\kappa_n+B_0[g(1)-g(3)],\qquad
 H_n(2)\le-\kappa_n+B_0g(2),
\]
where $B_0=2^{1-n}\binom nr$ and $B_0/\kappa_n=2n/(n+1)$.
The elementary bounds $7<e^2<37/5$ imply
\[
 g(1)>\frac{290}{441},\quad g(3)<\frac{14}{343},\quad
 g(1)-g(3)>\frac35,\quad g(2)<\frac{10}{49}<\frac12.
\]
As $5/3\le B_0/\kappa_n<2$, we obtain $H_n(1)>0$ and $H_n(2)<0$.
Strict decrease proves the unique global maximizing gain.

To check the affine Hessian, write $R=a_n\sqrt n$ and $q=q_n(a_n)>0$.
Permutation symmetry splits weight directions into $\operatorname{span}(\one)$
and its orthogonal complement. Bias--weight mixed derivatives vanish
under $X\mapsto-X$. The bias eigenvalue of $J_n$ is
\[
 \frac{D_n(a_n)}R-\frac q{R^3}<0.
\]
For a unit difference direction $(e_i-e_j)/\sqrt2$, conditioning on
$X_i=-X_j$ shows that the numerator's second derivative is
$D_{n-2}(a_n)$. Thus the common eigenvalue on $\one^\perp$ is
\[
 \frac{D_{n-2}(a_n)}R-\frac q{R^3}<0.
\]
Along the balanced line, at the stationary gain,
\[
 \frac{d^2}{da^2}\frac{q_n(a)}{a\sqrt n}\bigg|_{a=a_n}
 =\frac{q_n''(a_n)}{a_n\sqrt n}<0.
\]
Rescaling this direction to unit length preserves its sign. The gradient
vanishes by the stationary condition and symmetries. All Hessian
blocks are therefore negative definite.
\end{proof}

\begin{corollary}[A local weighted-code inequality]\label{cor:local-code}
Partition physical coordinates into $m$ disjoint pairs of rails
$(B_i^0,B_i^1)$, where $m$ is odd and each rail has odd size at least three.
For $s\in\F_2^m$ put $A_s=\bigcup_i B_i^{s_i}$ and $n_s=|A_s|$, and
assume $n_s\ge5$. Define
\[
 Q(b,u)=\frac{\sum_s w_{n_s}
      |\E[\chi_{A_s}\tanh(b+u\cdot X)]|}
                  {\sqrt{b^2+\|u\|_2^2}}.
\]
At every point with $b=0$, incoming weights $a_{n_s}$ on one selected
transversal $A_s$, and zero weights elsewhere, $Q$ has value one and a
strict local maximum. Input flips give the corresponding signed points.
\end{corollary}
\begin{proof}
The selected parity coefficient has the negative-definite quadratic
expansion in Theorem~\ref{thm:local}. Every inactive coordinate has the
same numerator second derivative as the bias, and mixed inactive
coordinates vanish by independence. Hence their normalized eigenvalue
is also $D_n/R-q/R^3<0$.

Each other transversal contains at least three inactive coordinates.
Repeated finite differences factor its coefficient by the product of
those incoming weights, so its absolute coefficient is
$O(\|(\Delta b,\Delta u)\|^3)$ near the selected point. The finite sum
of these terms does not change the quadratic part. The selected weight
$w_n$ makes its value exactly one. We use this second-order expansion;
the absolute-value objective need not be twice differentiable throughout
a neighborhood.
\end{proof}

Corollary~\ref{cor:local-code} does not assert $Q\le1$ globally. A
violating neuron could be far from every selected balanced atom.

\section{The sharp global bound for three-bit parity}

For three bits the global problem admits an exact algebraic solution.
Write
\[
 C_3(u)=\E[X_1X_2X_3\tanh(u\cdot X)],\qquad
 q_3(a)=\frac{2\tanh^3a}{1+3\tanh^2a}.
\]

\begin{theorem}[Global three-bit parity gain]\label{thm:three}
There is a unique maximizer $a_3>0$ of $q_3(a)/a$. Put
\[
 L_3=\frac{q_3(a_3)}{\sqrt3\,a_3}.
\]
Every affine tanh neuron, including neurons with additional irrelevant
input coordinates, satisfies
\[
 |\widehat{\tanh(u\cdot\bar X)}(\{1,2,3\})|
 \le L_3\|u\|_2.
\]
Nonzero equality requires zero bias, zero irrelevant weights, and equal
absolute relevant weights $a_3$. Consequently the variation cost of
three-bit parity is exactly $L_3^{-1}$.
\end{theorem}

\begin{proof}
First assume positive relevant weights $a,b,c$ and no bias. Let
$t_1=\tanh a,t_2=\tanh b,t_3=\tanh c$ and $s_i=t_i^2$. Direct use of
the addition formula gives
\[
 |C_3(a,b,c)|=\frac{2t_1t_2t_3\prod_i(1-s_i)}P,
\]
where
\[
 P=1-2\sum_{i<j}s_is_j+\sum_{i<j}s_i^2s_j^2
        -2s_1s_2s_3\sum_i s_i+8s_1s_2s_3.
\]
The denominator obeys the identity
\[
 P-\prod_i(1-s_i)\left(1+\sum_i s_i\right)
 =\frac12\sum_{\rm cyc}(s_i-s_j)^2(1-s_k)^2\ge0.
\]
Thus, by the arithmetic--geometric mean inequality,
\begin{equation}\label{eq:holder3}
 |C_3(a,b,c)|
 \le\frac{2t_1t_2t_3}{1+s_1+s_2+s_3}
 \le\prod_{i=1}^3\left(\frac{2t_i^3}{1+3t_i^2}\right)^{1/3}
 =\prod_i q_3(u_i)^{1/3}.
\end{equation}
Equality with positive weights requires all three equal. Input flips
handle negative weights, and a zero relevant weight makes the
coefficient zero. Applying $q_3(a)\le a\sup_t q_3(t)/t$ and then
geometric mean at most quadratic mean proves the sharp bias-free bound.

For uniqueness of the scalar gain, set $t=\tanh a$. The sign of
$(q_3(a)/a)'$ is the sign of
\[
 H(t)=3\operatorname{arctanh}(t)(1-t^4)-t(1+3t^2).
\]
Here $H(0)=0$, $H'(0)=2$, $H(1^-)=-4$, and
\[
 H'(t)=2-6t^2-12t^3\operatorname{arctanh}t,\qquad H''(t)<0
 \quad(0<t<1).
\]
Strict concavity gives exactly one zero in $(0,1)$, proving uniqueness.

We record the bias argument generally. For any odd activation and odd
$n$, its top parity coefficient satisfies
\begin{equation}\label{eq:bias-identity}
 C_n(B;u)=\frac12\left[
 C_n(0;u_1,\ldots,u_{n-1},u_n+B)
 +C_n(0;u_1,\ldots,u_{n-1},u_n-B)\right].
\end{equation}
Indeed, averaging the first $n-1$ coordinates with their parity gives
an odd function of the remaining scalar offset. If the bias-free gain
is at most $L\|u\|$, triangle inequality and root mean square yield
\[
 |C_n(B;u)|\le\frac L2\left[
 \sqrt{A+(u_n+B)^2}+\sqrt{A+(u_n-B)^2}\right]
 \le L\sqrt{A+u_n^2+B^2},
\]
where $A=\sum_{i<n}u_i^2$. For a nonzero coefficient, $u_n\ne0$,
so equality in root mean square forces $B=0$. Conditioning on irrelevant
coordinates and applying the same argument puts their squared norm
into the bias second moment. A nonzero irrelevant weight makes the
effective bias nonzero with positive probability and prevents equality.

Finally, signed averaging over the eight input flips projects a balanced
maximizing neuron to $q_3(a_3)\chi_{\{1,2,3\}}$, up to its fixed sign.
This gives a finite representation of parity with cost
$\sqrt3a_3/q_3(a_3)$, matching the lower bound from the sharp gain.
\end{proof}

Numerically $a_3\approx0.9655523$ and $L_3^{-1}\approx5.3674918$.
The numerical values are not used in the proof. The eventual
high-degree global theorem below uses a different argument; it does
not require a higher-degree version of \eqref{eq:holder3}.

\section{Positive curvature of the balanced parity cost}

\begin{theorem}[Optimized cost expansion]\label{thm:curvature}
Along odd integers $n\to\infty$, the balanced costs satisfy
\begin{equation}\label{eq:cost-asymptotic}
 w_n=An+B+\frac Cn+O(n^{-2}),\qquad C>0.
\end{equation}
They admit a holomorphic interpolation on
$\{\Re z>0,\ |z|>z_0\}$ with this uniform expansion. In particular,
\[
 w''(n)=\frac{2C}{n^3}+O(n^{-4}).
\]
Numerically $A\approx2.24178537487$, $B\approx-1.52970827506$, and
$C\approx0.59356404999$.
\end{theorem}

\begin{proof}
Write $J_j=j(j+1)$ and
\[
 R_j(z)=\prod_{\ell=0}^{j-1}\frac{z-(2\ell+1)}{z+(2\ell+3)},
 \qquad f_j(a)=\frac1{1+e^{2a(2j+1)}}.
\]
Pairing Hamming levels in \eqref{eq:central-sum} gives the exact identity
\begin{equation}\label{eq:Finterpolation}
 q_n(a)=\kappa_n F(n,a),\qquad
 F(z,a)=1-\frac{4z}{z+1}\sum_{j\ge0}(-1)^jR_j(z)f_j(a).
\end{equation}
At odd $n$, $R_j(n)=0$ for $j>(n-1)/2$. In the right half-plane each
factor of $R_j$ has modulus at most one, since
\[
 |z+(2\ell+3)|^2-|z-(2\ell+1)|^2
 =8(\ell+1)(\Re z+1)>0.
\]
The series is therefore locally normally convergent for $\Re a>0$.
Its large-$z$ expansion is
\[
 F(z,a)=F_0(a)+F_1(a)/z+F_2(a)/z^2+O(|z|^{-3}),
\]
where
\begin{align*}
 F_0&=1-4\sum_{j\ge0}(-1)^jf_j,\\
 F_1&=4\sum_{j\ge0}(-1)^j(1+2J_j)f_j,\\
 F_2&=-4\sum_{j\ge0}(-1)^j(1+4J_j+2J_j^2)f_j.
\end{align*}
For completeness, expand
\[
 R_j(z)=1-2J_j/z+(2J_j+2J_j^2)/z^2+O_j(|z|^{-3}).
\]
To sum the remainder uniformly, split at $j=\eta|z|$ for a sufficiently
small fixed $\eta$. Below the cutoff, the logarithmic Taylor remainder
is a polynomial in $j$ times $|z|^{-3}e^{O(j^2/|z|)}$; its exponential
is absorbed by $f_j$'s $e^{-4\Re(a)j}$ decay. Above the cutoff,
$|R_j|\le1$ makes the tail exponentially small. This argument holds
uniformly near the real interval $[1.237,1.238]$, including any fixed
number of $a$-derivatives.

The gamma-ratio interpolation of $\kappa_n$ has uniform expansion
\[
 \kappa_z=\sqrt{\frac2{\pi z}}
 \left(1+\frac1{4z}+\frac1{32z^2}+O(|z|^{-3})\right)
 \qquad(|\arg z|\le\pi/2).
\]
Set $D_1=F_1+F_0/4$ and $D_2=F_2+F_1/4+F_0/32$. Then
\[
 \frac{a\sqrt z}{\kappa_z F(z,a)}
 =A(a)z+B(a)+C_{\rm raw}(a)/z+O(|z|^{-2}),
\]
with
\[
 A(a)=\sqrt{\pi/2}\,\frac a{F_0},\qquad
 B(a)=-A(a)D_1/F_0,
\]
\[
 C_{\rm raw}(a)=A(a)\left[(D_1/F_0)^2-D_2/F_0\right].
\]
Let $a_0$ be the zero in $[1.237,1.238]$ of $aF_0'-F_0$.
The exact certificate below proves $F_0''<0$ there, so $A''(a_0)>0$.
The analytic implicit-function theorem gives a critical branch $a(z)$
near $a_0$ for all sufficiently large $z$ in the right half-plane.
For real odd $n$, Theorem~\ref{thm:local}'s uniqueness identifies it with
$a_n$. Substitution gives
\begin{equation}\label{eq:Cformula}
 A=A(a_0),\quad B=B(a_0),\quad
 C=C_{\rm raw}(a_0)-\frac{B'(a_0)^2}{2A''(a_0)}.
\end{equation}
This proves the claimed analytic remainder. Cauchy's estimate on a
radius-$n/2$ disk gives the second-derivative expansion.

Here are reproducible exact bounds for the sign in
\eqref{eq:Cformula}. At $a_0$, using $F_0'=F_0/a_0$,
\begin{equation}\label{eq:intervalC}
 \frac C{\sqrt{\pi/2}}
 =\frac{aD_1^2}{F_0^3}-\frac{aD_2}{F_0^2}
  +\frac{(D_1-aD_1')^2}{2aF_0''F_0^2}.
\end{equation}
The companion script \texttt{verification/certify\_curvature.py} uses
only rational arithmetic in its assertions. It encloses exponentials
by their Taylor sums through degree 256 and a geometric bound on the
remaining positive terms. It sums the six series
$F_0,F_1,F_2,F_0',F_0'',F_1'$ through $j=9$. On
$[1237/1000,619/500]$, each omitted tail is bounded by
\[
 192\cdot11^6 e^{-2(1237/1000)21}<10^{-14}.
\]
Indeed, every polynomial factor is bounded by $16(j+1)^6$, with the
common factor four, and a geometric ratio at most
$64e^{-4(1237/1000)}<64/100$; the sum of this geometric series is less
than three. Outward interval rounding uses denominator $10^{45}$.
The resulting rational enclosures prove opposite signs of
$aF_0'-F_0$ at the endpoints, $F_0''<0$ throughout the interval, and
\[
 0.4616< C/\sqrt{\pi/2}<0.4842.
\]
Thus $C>0$. Floating-point formatting is used only to display the
certificate's already proved interval endpoints.
\end{proof}

\input{source/code-atom.tex}

\input{source/overshoot-section.tex}

\input{source/cycle-ordering.tex}

\section{Completion of the online theorem}

The cycle construction of Section~\ref{sec:cycles} uses even rails and
one shared carrier. Its supports have odd degrees $n_s=\Theta(m^3)$.
The sharp global code-atom theorem applies to every affine neuron,
and its support bound is polynomial in these parameters.
Appendix~\ref{app:physical} therefore classifies every minimum-cost
completion after every calibrated history. The compulsory prefix
factors make the remaining Fourier costs small enough to absorb all
polynomial stability and rounding constants. The strict ordering in
Theorem~\ref{thm:cycle-order} supplies
$\exp(\Omega(\sqrt m))$ designated rounds.
Appendix~\ref{app:posterior} proves the conclusion for the actual
Gaussian posterior at all sufficiently large fixed widths.

It remains to verify the teacher resources. The target is the sign of
\[
 b(x)\left(\sum_{i=1}^m(-1)^{z_i(x)}-(m-1)\right)
 =\sum_{i=1}^m\chi_{A_{e_i}}(x)-(m-1)\chi_{A_0}(x),
\]
a score of margin one. Each parity has degree $O(m^3)$.
On the attainable values of the coordinate sum, parity is represented
by an alternating sum of threshold steps. Replace each step at a
midpoint by a tanh transition of gain $O(\log(mn_{\max}))$.
Each attainable value is at distance at least one from every midpoint,
so the uniform error can be made less than $1/(8m)$ per parity using
$O(n_{\max})$ neurons. The combined score has margin greater than
$1/2$; doubling its readout gives margin one. Its width is
$O(mn_{\max})=O(m^4)$, and incoming weights, biases, and readouts
are bounded by that width after increasing its fixed constant.
For the Boolean circuit, compute each rail parity and the base parity
once. Forming the logical bits uses only $O(\ell)$ further XOR gates,
since every cycle column has at most two nonzero entries. An AND of
the $m$ logical zero tests and a final sign multiplication then give
the target. The total bounded-fan-in gate count is $O(d_m)$.
This proves Theorem~\ref{thm:main-online}.

\appendix
\input{source/conditional-physical.tex}
\input{source/conditional-posterior.tex}
\input{source/periodic-appendix.tex}
\input{source/localization-appendix.tex}
\input{source/residual-clt.tex}
\input{source/equal-bulk-bias.tex}
\input{source/single-parity-appendix.tex}
\input{source/code-localization.tex}
\end{document}
